\scanpage{4}
\mspage{1}

\uline{Dérivateurs}

1. Préliminaires sur les foncteurs \unsure{image|images} \unsure{inverse|inverses} des préfaisceaux sur des catégories d'indices, et leurs adjoints multiples.

$A$ ``catégorie de coefficients''

$I$ ``catégorie d'indices'',

\[
A(I)=\underline{\mathrm{Hom}}(I^{\unsure{\circ|o}},A)\qquad\text{préf. sur }I\ \text{à valeurs }A
\]

Si $f:I\longrightarrow J$, on en déduit

\[
f^{*}:A(J)\longrightarrow A(I)
\]

\deleted{\unsure{E}}Et on s'intéresse à l'existence des foncteurs adjoints successifs à gauche et à droite. Notations provisoires

\[
f^{?}\ f_{!}\ f^{*}\ f_{*}\ f^{!}
\]

Voici les faits principaux :

a) Sous réserve de l'existence des $\varprojlim$ (resp $\underset{\unsure{-|\to}}{\lim}$) pertinentes, $f_{*}$ (resp $f_{!}$) existe. Par exemple s'il existe $\varprojlim$ (resp. $\varinjlim$) \uline{finies}, alors $f_{*}$ (resp. $f_{!}$) existe si $I,J$ sont des catégories finies.

b) Même si $A$ est une catégorie abélienne, \unsure{avec} $\varprojlim$ et $\varinjlim$\textsuperscript{*},\margin{* p. ex. \unsure{(Ab)}} les \uline{bi-adjoints} n'existent pas toujours\unsure{,|.} \deleted{\unsure{cf.}} \deleted{\unsure{déj.}} dans le cas particulier où $J=e$ est la cat. ponctuelle, et $I$ l'ens ordonné

\diagram{\begin{tikzcd}[ampersand replacement=\&]
\& a \\ \unsure{0|o}\arrow[ur]\arrow[dr] \& \\ \& b
\end{tikzcd}}

(dans le cas de $f_{\deleted{*}}^{!}$), ou l'ens ordonné opposé (dans le cas de $f^{?}$). Même si \deleted{c) Par contre, si} $f$ est une immersion \unsure{ouvert|ouverte}, $f^{!}$ n'existe pas forcément (par ex. \inserted{\unsure{définie par|déterminée par}} si $I$ comme dessus, et

\diagram{\begin{tikzcd}[ampersand replacement=\&]
\& a \arrow[dr] \& \\ J={}0\arrow[ur]\arrow[dr] \& \& c \\ \& b\arrow[ur] \&
\end{tikzcd}}

), et si $f$ est une imm. fermée, $f^{?}$ n'existe pas forcément.

c) Si $f$ est une \uline{imm. fermée}, et $A$ est \uline{pointée} (i.e. avec un objet à la fois initial et final) alors $f^{!}$ existe. Si $f$ est une imm. ouverte, $f^{?}$ existe. (Tjrs sous réserve de l'existence \inserted{dans $A$} des $\varprojlim$ites pertinentes.) \inserted{Si $f$ \unsure{une} immersion, $f_{*}$ et $f_{!}$ \unsure{plt fid|pln fid} i.e. $f^{*}$ \unsure{un foncteur localisation}.}

d) Supposons seulement l'existence de l'objet final\unsure{\textsuperscript{*}} dans $A$ (\unsure{plus|pas} \unsure{on|n'a} besoin autres $\varprojlim$ finies). Si \deleted{\unsure{$f$ immersion fermée}}

% archive sheet number: 3
